\section{工作组织}
\label{sec:work_organization}


下面给出使用 \tkzNamePack{tkz-euclide} 与 \tkzEngine{Lua\LATEX} 时的一种示例组织方式。

\begin{itemize}
  \item \texttt{编译方式：} 添加如下一行：
  \begin{verbatim}
  % !TEX TS-program = lualatex
  \end{verbatim}
  以确保文档使用 \tkzEngine{Lua\LATEX} 编译。

  \item \texttt{文档类：} 如果目的只是绘制一幅图形，推荐使用 \texttt{standalone} 文档类，因为它可避免不必要的开销。

  \item \texttt{宏包加载：} 你可以用两种方式加载 \tkzNamePack{tkz-euclide}：
  \begin{itemize}
    \item \tkzcname{usepackage\{tkz-euclide\} } 让你获得整个宏包的完整功能。
    \item 推荐的做法是使用 \tkzPopt{tkz-euclide}{mini} 选项，它只加载绘图所需的模块。如果你愿意，仍然可以直接用 \TIKZ{} 绘图。
  \end{itemize}

  \item \texttt{条件判断：} 当你需要在文档中求值布尔条件时，\tkzNamePack{ifthen} 宏包很有用。

  \item \texttt{Lua 代码组织：} 虽然你可以直接用 \verb|\directlua| 嵌入 Lua 代码，但把代码外置有如下几个好处：
  \begin{itemize}
    \item 在支持 \tkzname{Lua} 和 \LATEX{} 的编辑器中，可以获得更好的语法高亮与代码展示效果。
    \item 注释更简单：\tkzname{Lua} 使用 \verb|--|，而 \LATEX{} 使用 \verb|%|。把 Lua 放在单独的文件里可避免混淆。
    \item 可复用：外部代码文件可在多个文档或图形之间复用。
  \end{itemize}
\end{itemize}

为简单起见，本文档在大多数情况下使用内嵌的 Lua 代码。不过在某些示例中，会使用外部文件来向你展示具体做法。


\vspace{12pt}
\begin{mybox}{}
\begin{minipage}{.50\textwidth}
\begin{verbatim}
% !TEX TS-program = lualatex
% Created by Alain Matthes on 2024-01-09.
\documentclass[margin = 12pt]{standalone}
\usepackage[mini]{tkz-euclide}
\usepackage{tkz-elements,ifthen}

\begin{document}
\directlua{
 init_elements() % Clear tables
 dofile ("lua/sangaku.lua")
 % Load the lua code
}
\begin{tikzpicture}[ scale = .75]
 \tkzGetNodes
 \tkzDrawCircle(I,F)
 \tkzFillPolygon[purple](A,C,D)
 \tkzFillPolygon[blue!50!black](A,B,C)
 \tkzFillCircle[orange](I,F)
\end{tikzpicture}
\end{document}
\end{verbatim}
\end{minipage}
\begin{minipage}{.40\textwidth}
\directlua{
 init_elements()
 dofile ("lua/sangaku.lua")
}
\begin{center}
\begin{tikzpicture}[ scale = .75]
  \tkzGetNodes
  \tkzDrawCircle(I,F)
  \tkzFillPolygon[purple](A,C,D)
  \tkzFillPolygon[blue!50!black](A,B,C)
  \tkzFillCircle[orange](I,F)
\end{tikzpicture}
\end{center}
\end{minipage}
\end{mybox}

\vspace{12pt}
下面是 \tkzname{Lua} 部分的代码：文件 |sangaku.lua|

\vspace{1em}
\begin{minipage}{.5\textwidth}
\begin{mybox}
\begin{verbatim}
 z.A = point(0, 0)
 z.B = point(8, 0)
 L.AB = line(z.A, z.B)
 S.AB = L.AB:square()
 _, _, z.C, z.D = S.AB:get()
 z.F = S.AB.ac:projection(z.B)
 L.BF = line(z.B, z.F)
 T.ABC = triangle(z.A, z.B, z.C)
 L.bi = T.ABC:bisector(2)
 z.c = L.bi.pb
 L.Cc = line(z.C, z.c)
 z.I = intersection(L.Cc, L.BF)
\end{verbatim}
\end{mybox}
\end{minipage}


\subsection{缩放策略的更新}

在以前的版本中，建议在代码的 Lua 部分进行缩放。不过，这一指导原则现在已经改变。

由于所有几何计算都在 \tkzname{Lua} 中完成，在 \TIKZ{} 中进行缩放已不再有任何问题。相反，在 \tkzname{Lua} 中进行缩放却带来了若干麻烦——尤其是最近实现的与圆锥曲线相关的函数，它们涉及大量使用实数的距离计算。

这些问题促使我们重新审视了若干函数，并在此过程中发现并修复了一些与 Lua 端缩放相关的 bug。

\medskip

\textbf{\textcolor{red}{新的建议：}} 从现在起，缩放应只在 \TIKZ{} 部分进行。\tkzname{Lua} 代码应在一个未经缩放、一致的坐标系中运行，以确保所有几何计算的可靠性。

本文档后续只在 \tkzEnv{tikz}{tikzpicture} 环境中使用缩放。
\endinput